\documentclass[11pt]{article}
\usepackage{amsmath,amssymb}
\usepackage[margin=1in]{geometry}
\usepackage{microtype}
\setlength{\parindent}{1.5em}
\allowdisplaybreaks

\title{Rewritten Chain of Thought: Refuting the Single-Rational-Function Hint
and Reaching Only a Chamber-Restricted $A_n$\\
in the Two-Minus Sector}
\author{}
\date{}

\begin{document}
\maketitle

\begin{abstract}
\noindent This document is a first-person reconstruction of the chain of thought for an
AI attempt to find a closed-form tree-level amplitude $A_n$ in the two-minus sector of
1D deep-water surface waves, working only from the prompt, the code \texttt{OnShellBG.m},
and the data that code produces. It is rebuilt from the visible run trajectory --- the
raw private chain of thought was encrypted, so this is reconstructed from the visible
messages, tool calls, and outputs --- and it keeps the hypotheses I formed, the concrete
kinematic points I tested, the numbers the Berends--Giele code returned, and the
dead-ends, in the order they actually happened. I should be honest about the outcome
up front. The prompt's hint insisted that $A_n$ is a \emph{single global rational
function} with simple poles on factorization channels, no chambers, no absolute values,
and ``emphatically not a plain polynomial,'' and it told me to fit one $N/D$ ansatz on
generic kinematics. I did not confirm that picture; I refuted it. The data forced a
genuinely \emph{piecewise} amplitude (the $|k_S|$ in the propagators flip sign), and I
proved that no single global rational function can exist. What I delivered is therefore a
\emph{chamber-restricted} answer: the clean monomial
$A_n=2^{\,n-1}i\,\omega_1\,\omega_2^{\,2n-5}$, exact (verified to relative error $0$)
only in the chamber selected by the benchmark's standard ascending-positive sampling,
together with its equivalent free-frequency ratio form and an explicit statement that the
fully general object is an irreducibly $\min/\max$, sorted-magnitude, piecewise spline
that I did not fold into one expression. So this is a partial result that saw through the
hint rather than the all-kinematics closed form the task asked for.
\end{abstract}

\section*{Rewritten Chain of Thought}

\noindent\textbf{Orienting myself.}
The object is the on-shell tree amplitude $A_n$ in the sector
$\sigma=(-1,-1,+1,\dots,+1)$: legs $1,2$ are ``minus'', legs $3,\dots,n$ are ``plus''.
The dispersion ties momentum to frequency, $\omega_i^2 = g|k_i|$ with
$k_i = \sigma_i\omega_i^2/g$, and the resonant manifold imposes $\sum_i\omega_i=0$ and
$\sum_i\sigma_i\omega_i^2=0$. The recursion glues vertices with a propagator
\[
\frac{-i}{\;\omega_S^2/|k_S|-g\;},\qquad
k_S=\sum_{i\in S}k_i=\tfrac1g\sum_{i\in S}\sigma_i\omega_i^2 .
\]
The hint wants $A_n=N(\omega)/D(\omega)$ with $D$ a product of channel factors, fit by
matching at many points, a \emph{single} global rational function with no $|\cdot|$, no
$\min/\max$, no chambers, and not a plain polynomial. I keep that target in mind, but the
$|k_S|$ in the propagator is an immediate red flag, so I resolve to test the global-rational
claim rather than assume it.

\medskip
\noindent\textbf{First contact: amplitudes and the scaling weight.}
My first script just runs the code. At one generic point per $n$ I get
\[
A_5=-\tfrac{2349}{4}i,\qquad A_6=-\tfrac{483327}{152}i,\qquad A_7=-\tfrac{7302393}{400}i,
\]
with $A_7$ already costing about $19$ seconds in exact arithmetic --- a hint that I must be
economical at high $n$. My $n=4$ point returns \texttt{Indeterminate}; I note it and move on,
suspecting an accidental degeneracy ($k_{13}=0$ at the point I chose). I then scale all
frequencies $\omega\to\lambda\omega$ and read off the weight: $A_5$ multiplies by $2^6$,
$A_6$ by $2^8$. So
\[
\deg A_n = 2n-4 .
\]
Hypothesis 1, the bookkeeping one: $A_n=N/D$ is homogeneous of degree $2n-4$ (counting
$\deg\omega_i=1$, since $k\sim\omega^2$ and the propagator is weight $0$).

\medskip
\noindent\textbf{Thinking about the denominator from physics.}
Where are the poles? A propagator goes on shell at $\omega_S^2=|P_S|$, where
$P_S\equiv\sum_{i\in S}\sigma_i\omega_i^2$. External legs are safe, $|k_i|=\omega_i^2/g$,
because $\omega_i^2\ge0$. Internal lines are the worry. Splitting subsets by how they meet
$\{1,2\}$:
\begin{itemize}
\item all-plus $S\subseteq\{3,\dots,n\}$: $P_S=\sum_{i\in S}\omega_i^2>0$, definite, and the
channel factor is the clean polynomial $\omega_S^2-\sum_{i\in S}\omega_i^2=2e_2(S)$ with
$e_2(S)=\sum_{i<j\in S}\omega_i\omega_j$;
\item $S$ meeting exactly one of $\{1,2\}$: $P_S=-\omega_a^2+(\text{plus})$, sign indefinite.
\end{itemize}
Hypothesis 2: maybe $D_n=\prod_{S\subseteq\{3..n\},|S|\ge2}2e_2(S)$ (the safe channels only),
and the mixed channels carry no pole. For $n=5$ that would be
$D_5=(\omega_3\omega_4)(\omega_3\omega_5)(\omega_4\omega_5)(\omega_3\omega_4+\omega_3\omega_5+\omega_4\omega_5)$,
degree $8$, forcing $\deg N=2n-4+8=14$ --- uncomfortably large, and it presupposes the
mixed channels are inert. I distrust this enough to test the rationality claim directly
before committing to any $D$.

\medskip
\noindent\textbf{The bottom of the tower forces my hand ($n=4$).}
At $n=4$ the two conservation laws on four frequencies are rigid. Taking $\omega_2,\omega_3$
free and solving, the solution is \emph{forced}:
\[
\omega_1=-\omega_3,\qquad \omega_4=-\omega_2,
\]
a two-parameter family $(-\omega_3,\omega_2,\omega_3,-\omega_2)$. (This also explains my
earlier \texttt{Indeterminate}: with $\omega_1=-\omega_3$ one gets $k_2+k_4=0$ and
$\omega_2+\omega_4=0$, so a sub-current propagator is $0/0$.) Computing the amplitude
symbolically and simplifying with the magnitudes kept symbolic gives a \emph{piecewise}
answer:
\[
A_4=-8i\,\omega_2^3\omega_3\ \ (\omega_2<\omega_3),\qquad
A_4=-8i\,\omega_2\omega_3^3\ \ (\omega_2>\omega_3).
\]
This is decisive and unwelcome: in signed variables the pieces are
$-8i\,\omega_1\omega_4^3$ and $-8i\,\omega_1^3\omega_4$, switching across
$\omega_2^2=\omega_3^2$. By the identity theorem a single rational function equal to
$-8i\,\omega_2^3\omega_3$ on an open set must equal it everywhere; it does not. So $A_4$ is
genuinely piecewise. Hypothesis 3 (\emph{the amplitude is globally rational}) is in serious
doubt. I conjecture instead that the hint's ``single rational function'' holds only
chamber-by-chamber, and that the benchmark's sampling lives in one chamber.

\medskip
\noindent\textbf{A method to see chambers: resolve $|\cdot|$ by a reference sign.}
To get an honest rational function I evaluate every $|x|$ in the symbolic amplitude by the
sign of $x$ at a chosen numeric reference point; that is exactly valid throughout the open
sign-chamber containing the reference, and I can then \texttt{Together}/\texttt{Factor}. At
$n=5$ the raw amplitude has leafcount $\sim4\times10^5$, but the chamber with free
frequencies $(a,b,c)=(\omega_2,\omega_3,\omega_4)=(2,3,5)$ collapses to
\[
A_5=\frac{-16\,i\,a^5\bigl(ab+b^2+ac+bc+c^2\bigr)}{a+b+c}.
\]
Two other reference chambers, $(7,11,2)$ and $(3,5,-2)$, give visibly \emph{different}
rational functions. Hypothesis 3 is now refuted outright: $A_5$ is chamber-dependent.

\medskip
\noindent\textbf{Mapping where this chamber holds.}
I test $f_A(a,b,c)=-16i\,a^5(ab+b^2+ac+bc+c^2)/(a+b+c)$ against the recursion. It matches at
$(2,3,5)$ and $(2,5,3)$ but not at $(3,2,5)$ (ratio $56/81$), $(5,2,3)$ (ratio $72/625$),
$(3,5,2)$, $(5,3,2)$. Reading the pattern: $f_A$ is exact precisely when $a=\omega_2$ is the
\emph{smallest} of $a,b,c$, and it is symmetric in $b\leftrightarrow c$ (i.e.\ in the two
free plus legs). So the chamber is ``$\omega_2$ smallest''.

\medskip
\noindent\textbf{The collapse to a monomial.}
The denominator $a+b+c$ is exactly the quantity that blows up in the kinematic solver, so it
should be a coordinate artifact. Using the conservation laws I find the identities
\[
ab+b^2+ac+bc+c^2=\omega_1(\omega_1+\omega_5),\qquad a+b+c=-(\omega_1+\omega_5),
\]
and the artifact cancels:
\[
A_5\big|_{\omega_2\ \mathrm{smallest}}=16\,i\,\omega_1\,\omega_2^5 .
\]
A pure monomial. I check it head-on: $(2,3,5)\!\mapsto\!-3328i$, $(3,5,7)\!\mapsto\!-37584i$,
$(1,4,9)\!\mapsto\!-1168i/7$ --- all exact. As a function it uses only the two minus-leg
frequencies, with the high power on the smaller one. This is beautiful but it is a
\emph{polynomial}, contradicting ``not a plain polynomial'': another sign that the hint
describes a per-chamber rational form, not the on-shell value.

\medskip
\noindent\textbf{Cataloguing the other chambers ($n=5$).}
I grind out five more reference chambers and reduce each to $\omega$ variables:
\begin{align*}
\omega_2\ \text{smallest}:&\quad 16i\,\omega_1\omega_2^5,\\
\omega_3\ \text{smallest},\ \omega_2\ \text{next}:&\quad 16i\,\omega_1\omega_2\,\omega_3^2(2\omega_2^2-\omega_3^2),\\
\omega_3,\omega_4\ \text{smallest}:&\quad 32i\,\omega_1\omega_2\,\omega_3^2\omega_4^2,\\
\omega_5,\omega_3\ \text{smallest}:&\quad 32i\,\omega_1\omega_2\,\omega_3^2\omega_5^2,\\
\omega_5,\omega_2\ \text{smallest}:&\quad 16i\,\omega_1\omega_2\,\omega_5^2(2\omega_2^2-\omega_5^2).
\end{align*}
(For the last two I needed an auxiliary identity: the quartic $Q$ appearing there equals
$a^2s^2-2abcs-b^2c^2$ with $s=a+b+c$, and $Q/s^2=2\omega_2^2-\omega_5^2$ on shell.) Every
chamber reduces to a \emph{polynomial}; the solver artifact always cancels. So $A_5$ is a
symmetric, piecewise-\emph{polynomial} function whose pieces are labelled by the ordering of
the $|\omega_i|$ --- precisely the $\min/\max$ structure the hint forbids. Hypothesis 4
(\emph{some chamber has a genuine non-polynomial pole}) fails too: on shell there are no
poles at all.

\medskip
\noindent\textbf{Symmetry check.}
$16i\,\omega_1\omega_2^5$ looks asymmetric under $1\leftrightarrow2$, but the recursion picks
leg $1$ as a root, so I test directly: permuting the $(k,\omega)$ pairs --- swap the minus
legs, swap two plus legs, cycle $3\!\to\!4\!\to\!5$, or an arbitrary permutation
$\{3,1,5,2,4\}$ --- all return $-3328i$ at the $(2,3,5)$ point. So $A_n$ is fully symmetric
under same-$\sigma$ relabelling; the monomial is a chamber \emph{representative}, valid when
$\omega_2$ happens to be the smallest magnitude.

\medskip
\noindent\textbf{The conjecture and its test ($n=4,5,6,7$).}
The clean chamber generalises. The coefficients $8,16,\dots$ at $n=4,5$ suggest $2^{n-1}$ and
the power $2n-5$, so
\[
\boxed{\,A_n=2^{\,n-1}\,i\,\omega_1\,\omega_2^{\,2n-5}\,}\qquad\text{when }|\omega_2|=\min_i|\omega_i| .
\]
I throw a spread of points at it. At $n=5$: $(2,3,5),(1,4,9)$ exact. At $n=6$:
$(2,3,5,7)\!\mapsto\!-\tfrac{753664}{17}i$, $(1,2,4,8)\!\mapsto\!-\tfrac{4928}{15}i$,
$(3,4,5,6)\!\mapsto\!-762048i$, and crucially the lopsided $(2,3,100,101)$ also matches ---
so it is not merely a ``comparable magnitude'' coincidence. At $n=7$:
$(2,3,5,7,11)\!\mapsto\!-\tfrac{4030464}{7}i$ and $(1,2,3,4,5)\!\mapsto\!-\tfrac{8896}{15}i$.
At $n=4$ the recursion is \texttt{Indeterminate} (the forced $0/0$), but the formula returns
$-192i$ for $\{2,3\}$ and $-40i$ for $\{1,5\}$, matching the symbolic $n=4$ limit
$8i\,\omega_1\omega_2^3=-8i\,\omega_2^3\omega_3$. Every test is exact.

\medskip
\noindent\textbf{Recovering the ``ratio with a pole'' the hint wants.}
If I keep the solved frequencies eliminated rather than substituted, the monomial becomes a
genuine rational function of the free frequencies. With $S_1=\omega_2+\dots+\omega_{n-1}=
-(\omega_1+\omega_n)$ and
$\omega_1=-(S_1^2-\omega_2^2+\sum_{i\ge3}\omega_i^2)/(2S_1)$,
\[
A_n=-\,2^{\,n-2}\,i\,
\frac{\omega_2^{\,2n-5}\bigl(S_1^2-\omega_2^2+\sum_{i=3}^{n-1}\omega_i^2\bigr)}{S_1}.
\]
For $n=5$ the bracket is $2(ab+b^2+ac+bc+c^2)$, reproducing the chamber form; for $n=6$ at
$(2,3,5,7)$ one has $S_1=17$ and bracket $=289-4+83=368$, giving $-16i\cdot128\cdot368/17=
-\tfrac{753664}{17}i$, as it must. The numerator does not vanish at $S_1=0$ (for $n=5$ it
reduces to $-\omega_3\omega_4$), so the pole on the channel sub-energy $S_1$ is genuine ---
this is the ``ratio of polynomials, simple channel-sub-energy pole, not a plain polynomial''
the hint describes, now correctly understood as a per-chamber statement.

\medskip
\noindent\textbf{Trying for the fully general (any-chamber) rule.}
For robustness I attempt the generic $n=5$ law. Writing $A_5=16i\,(\omega_1\omega_2)\,\Phi$
with $\Phi$ a degree-four factor, the chambers above suggest that $\Phi$ depends only on the
two smallest magnitudes $f_1\le f_2$ and their signs $s_1,s_2$ ($-1$ for a minus leg):
\[
\Phi=\begin{cases}
f_1^4,& s_1=-1,\\
f_1^2(2f_2^2-f_1^2),& s_1=+1,\ s_2=-1,\\
2f_1^2f_2^2,& s_1=+1,\ s_2=+1.
\end{cases}
\]
(I also note that both minus legs can never be the two smallest: the momentum law gives
$\kappa_1+\kappa_2=\kappa_3+\kappa_4+\kappa_5$ with $\kappa_i=\omega_i^2$, and if
$\kappa_1,\kappa_2$ were smallest then $\kappa_3+\kappa_4+\kappa_5\ge3\max(\kappa_1,\kappa_2)
>2\max\ge\kappa_1+\kappa_2$, a contradiction. So the ``$s_1=-1$'' case always has $s_2=+1$.)
Testing this rule at $40$ random points gives $38$ exact matches; the two failures are
instructive --- one point has $\omega_5=0$ (a genuine degeneracy, recursion
\texttt{Indeterminate}), and one has \emph{negative} free frequencies, where a
$(+,+)$-smallest configuration disagrees ($-4.927\times10^6 i$ from the code versus
$-5.690\times10^6 i$ from the rule). So once signs are mixed the chamber structure is finer
than ``the two smallest'', confirming that the genuinely global object is an irreducibly
$\min/\max$/sorted thing. The clean single formula is the canonical-chamber one; the general
case is the symmetric piecewise extension.

\medskip
\noindent\textbf{Why no single rational function can exist.}
I make the obstruction precise. The channel $\{2,3,4\}$ has
$P_{234}=-\kappa_2+\kappa_3+\kappa_4$, and using $\kappa_3+\kappa_4=\kappa_1+\kappa_2-\kappa_5$
this is $P_{234}=\kappa_1-\kappa_5$, whose sign is $\mathrm{sgn}(|\omega_1|-|\omega_5|)$ ---
not fixed across the sector. So the unregularised $|k_S|$ genuinely flips, and the recursion
is piecewise. Since $A_5$ equals the monomial $16i\,\omega_1\omega_2^5$ on the entire open
chamber $\{|\omega_2|=\min\}$, any rational function agreeing with the code there must
\emph{be} that monomial, which disagrees in the adjacent chamber. Hence no single global
rational function exists; the hint's ``one fit valid for all the sector'' is achievable only
within a chamber.

\medskip
\noindent\textbf{Choosing the chamber and verifying hard.}
The natural chamber is the one every example in the code inhabits: positive free frequencies
in ascending order, so that $\omega_2$ (the lone free minus leg, smallest of the free set)
ends up the global minimum magnitude. I confirm $\omega_2=\arg\min|\omega|$ holds at every
such point. Against the \emph{provided} \texttt{BGAmplitude} the boxed formula then matches
\emph{exactly} (relative error $0$): $30/30$ points at $n=5$, $18/18$ at $n=6$, $8/8$ at
$n=7$, plus explicit values like $\{3/2,2,5/2\}\!\mapsto\!-\tfrac{891}{2}i$ ($n=5$),
$\{1,3/2,2,5/2\}\!\mapsto\!-\tfrac{968}{7}i$ ($n=6$),
$\{1,3/2,2,5/2,3\}\!\mapsto\!-\tfrac{1928}{5}i$ ($n=7$).

\medskip
\noindent\textbf{An independent re-implementation as insurance.}
Because matching one code against a formula derived from it can hide a shared bug, I port the
Berends--Giele recursion to Python from scratch --- once in floating point and once in
\emph{exact} Gaussian-rational arithmetic (a small $a+bi$ class over the rationals). This
separate codebase reproduces the Wolfram reference values
($(2,3,5)\!\mapsto\!-3328i$, $(2,3,5,7)\!\mapsto\!-\tfrac{753664}{17}i$,
$(2,3,5,7,11)\!\mapsto\!-\tfrac{4030464}{7}i$) and agrees \emph{exactly} with the closed
form over a fresh random sweep: $40/40$ at $n=5$, $20/20$ at $n=6$, $6/6$ at $n=7$. I also
confirm the two forms of the answer (the all-frequency monomial and the free-frequency ratio)
agree at $50$ points each for $n=4,5,6,7$.

\medskip
\noindent\textbf{Closing $n=4$ honestly.}
The recursion cannot evaluate $n=4$ directly: the forced kinematics make $k_2+k_4=0$ and
$\omega_2+\omega_4=0$, so the $\{2,4\}$ propagator is a removable $0/0$ and the bare code
returns \texttt{Indeterminate}. Taking the finite limit with $|\cdot|$ symbolic gives
$A_4=-8i\,\omega_2^3\omega_3=8i\,\omega_1\omega_2^3$ for $\omega_2<\omega_3$ (and the swapped
$-8i\,\omega_2\omega_3^3$ otherwise), exactly the formula; numeric limits
$\{3/2,5/2\}\!\mapsto\!-\tfrac{135}{2}i$, $\{1,4\}\!\mapsto\!-32i$,
$\{2,7\}\!\mapsto\!-448i$ confirm it.

\medskip
\noindent\textbf{What the structure is telling me.}
The amplitude vanishing like $\omega_2^{\,2n-5}$ as the smallest frequency goes soft is a
strong zero, not the single power one might expect, and it is the same phenomenon behind the
$n=4$ result $-8i\,\omega_2^3\omega_3$. The mechanism is that conservation pins the two minus
legs to the extremes of the magnitude ordering with the plus legs in the middle, so that
everything except the softest minus leg is fixed and that leg carries the remaining mass
dimension by itself. The clean statement is the monomial $A_n=2^{n-1}i\,\omega_1\omega_2^{2n-5}$
(equivalently the displayed ratio with its $S_1$ pole), exact in the chamber selected by the
standard ascending sampling; the genuinely global object is its symmetric, sorted-magnitude,
piecewise extension, and no single rational function can represent it because the internal
$|k_S|$ honestly changes sign within the sector.

\end{document}
